To explore the proposed method, an illustrative problem is considered first.
The problem consists of a simulation of a simple 1D mass-spring-damper system (\Cref{fig:toy-setup}) and an imperfectly tuned PID controller.
The mass is suspended in a hanging position with gravity acting on it continuously.

\begin{figure}[htb]
    \centering
    \def\svgwidth{0.33\linewidth}
    \small\selectfont\input{Thesis/figures/method/toy_setup.pdf_tex}
    \caption{Setup of the mass-spring-damper system in the illustrative problem with positive directions indicated.}
    \label{fig:toy-setup}
\end{figure}

The system state $\mathbf{x}$ consists of the position and velocity $x,\,\dot{x}$, respectively.
The controller aims to reach targets of the form $\mathbf{x}_d = [x_d,\, 0]$, but controls position only.
At regular intervals, a random target position $x_d$ is generated for the controller to reach and maintain.
Each timestep, the adapter calculates the required $\Delta$, and the controller computes the control action based on the current position and the adapted target position.
The system's response to this control input is simulated and the tuple $(x_0, \dot{x}_0, u, x_d, 0)$ is added to the buffer.
This buffer stores a limited number of tuples corresponding to a specified window of the most recent data:
Once capacity is reached, the oldest element is removed at every timestep.
At regular update intervals, the adapter is trained on the data in the buffer.
To ensure meaningful training, it is important that the tuples in the training data always contain the \emph{true} target position $x_d$, and not the target that was actually provided to the controller, i.e.,~$x_d + \Delta$.
As the simulation progresses, the parameters of the dynamic system are gradually altered to simulate system degradation.
The pseudo algorithm for the implementation can be found in~\Cref{alg:toy-problem}.

\begin{algorithm}
    \SetAlgoLined
    \caption{Target adaptation}
    \label{alg:toy-problem}

    \KwData{system state $\mathbf{x}$, target $\mathbf{x}_d$, adaption $\Delta$, control action $u$}

    Initialize model:  \textit{adapter}\;
    Initialize buffer: $\textit{buffer} \gets []$\;
    \For{each timestep $t_i$}{
        \If{len(buffer) > threshold}{
            Remove the oldest element in the buffer\;
        }
        \If{$t_i \mod target\_interval = 0$}{
            $\mathbf{x}_d \gets random\_target()$\;
        }
        \If{$t_i \mod update\_interval = 0$}{
            $(\textit{features, labels}) \gets \textit{prep\_data(buffer, window)}$\;
            $\textit{adapter.fit(features, labels)}$\;
        }

        $\Delta \gets adapter.predict([\mathbf{x}_0, \mathbf{x}_d])$\;
        $u \gets controller.compute\_control(\mathbf{x}_0, \mathbf{x}_d + [\Delta, 0])$\;
        Simulate the system response $\mathbf{x}$\;
        Update system parameters\;
        Append $(\mathbf{x}_0, u, \mathbf{x}_d)$ to \textit{buffer}\;
        $\mathbf{x}_0 \gets \mathbf{x}$\;
    }
\end{algorithm}

The simulation ran for 60\si{\second} at 50\si{\hertz}.
The target was changed every 5\si{\second}, and the adapter was trained fully online, updating every 15\si{\second}.
Training features were created for multiple prediction windows, this means the training features consisted of pairs $[\mathbf{x}_t, \mathbf{x}_{t+\textit{window}}]$ for windows of 3, 5, 10, 15, and 25.
The buffer's budget corresponded to the past 30\si{\second} of data.
As a baseline, a second controller and system were also simulated, identical to the described simulation, but without the addition of an adapter.
To both the baseline and adapted signal, a small amount of Gaussian noise was added to simulate measurement noise.

\subsection{An Over-Tuned Controller}\label{subsec:toy-aggressive}
\begin{figure*}
    \centering
    \includestandalone[build={latexoptions={"-output-directory=./Thesis/figures/toy-problem/"}}]{Thesis/figures/toy-problem/toy_plot}
    \vspace*{-2em}
    \caption{The full minute simulation of the system with an aggressive controller. \emph{Top}: The system response $x$ and target $x_d+\Delta$ of the adapted controller, alongside the response and target of the unadapted PID controller ($\Delta=0$).
    \emph{Bottom}: The resulting adapted control inputs from the controller receiving $x_d+\Delta$ alongside the unaltered PID control inputs.
    During the first 15\si{\second} (3 targets) the adapter has not received updates yet, and the signal is thus identical to the baseline PID signal.}
    \label{fig:toy-plot}
\end{figure*}

For a first simulation, a combination of a system and controller was chosen where the controller's gains were tuned too high.
In~\Cref{fig:toy-plot}, the result of the simulation can be observed.
At the start, due to the near-zero initialization of the adapter's parameters, the signal response of the system controlled by the adapted controller is indistinguishable from that of the baseline PID controller.
However, as soon as the adapter has trained on the buffer, the changes to the target become visible.

The effects of the adaptation can be more clearly observed in~\Cref{fig:toy-plot-closeup}.
Since the baseline PID controller is too aggressive, the signal overshoots and oscillates significantly before settling.
In contrast, the adapter counters this by bringing the target closer to the current state.
This is the logical result of the buffer containing data corresponding to the overshoots, where the adapter would often encounter training features for which the labeled target position\footnote{Rather, the labeled $\Delta = x_{d} - x_{t+\emph{window}}$, but the explanation provided is equivalent and somewhat clearer.} $x_d$ is somewhere in between $x_t$ and $x_{t+\emph{window}}$.
Therefore, upon receiving the input feature $[\mathbf{x}_t, \mathbf{x}_d]$
%(representing $\mathbf{x}_d$ being reached from $\mathbf{x}_t$ within 10 timesteps),
it produces and adapted target position in between $x_t$ and $x_d$.
As the state approaches the adapted target, the adapter starts producing targets closer to the reference target, guiding the system towards this desired state.

\begin{figure}
    \centering
    \includestandalone[build={latexoptions={"-output-directory=./Thesis/figures/toy-problem/"}}]{Thesis/figures/toy-problem/toy_plot_closeup}
    \vspace*{-2em}
    \caption{A close-up of the second to last target in the one-minute simulation. The effects of the adapter can be more clearly observed: The overshoot (seen in the baseline PID signal) is countered by bringing the target closer to the current state.}
    \label{fig:toy-plot-closeup}
\end{figure}

\subsection{An Under-Tuned Controller}\label{subsec:toy-damped}
For a second simulation, the system and controller were designed such that the controller was more conservative or \emph{under-tuned}.
Additionally, the signal produced by the baseline PID controller has a clear steady-state error.
The result of the simulation can be found in~\Cref{fig:toy-plot-damped}.

\begin{figure*}
    \centering
    \includestandalone[build={latexoptions={"-output-directory=./Thesis/figures/toy-problem/"}}]{Thesis/figures/toy-problem/toy_plot_damped}
    \vspace*{-2em}
    \caption{The full minute simulation of the system with an under-responsive controller. \emph{Top}: The system response and target of the adapted controller, alongside the response and target of the unadapted baseline PID controller.
    \emph{Bottom}: The resulting adapted control inputs from the controller receiving the altered target alongside the unaltered baseline PID control inputs.}
    \label{fig:toy-plot-damped}
\end{figure*}

Once more, the adapter is able to improve the controller's performance after just a single update.
In~\Cref{fig:toy-plot-damped-closeup} it can be seen that the adapted target is adjusted away from the current state.
This is because training features created from the data in the buffer frequently include state pairs $\mathbf{x}_t,\, \mathbf{x}_{t+\emph{window}}$ where the target position $x_d$ is positioned further from $x_t$ than it is from $x_{t+\emph{window}}$.
As a result, the adapter learns to predict adaptations that would move $x_d$ farther from $x_t$, making the controller more aggressive.
Another effect of the adapter is the consistent offset applied to the target:
In the training features, the adapter regularly encounters pairs $\mathbf{x}_t, \, \mathbf{x}_{t+\emph{window}} = \mathbf{x}_t$, which appear to represent the target state being reached, even though $x_{t+\emph{window}}\neq x_d$.
These features are the result of the steady-state error present in the signal, leading the adapter to learn to apply an offset to the target to counteract this.

\begin{figure}
    \centering
    \includestandalone[build={latexoptions={"-output-directory=./Thesis/figures/toy-problem/"}}]{Thesis/figures/toy-problem/toy_plot_damped_closeup}
    \vspace*{-2em}
    \caption{A close-up of the second to last target in the one-minute simulation. The target is initially moved further from the current state to `pull' it towards the desired state in a shorter period. In addition, the adapted target is offset from the desired state as to counter the steady-state error seen in the baseline PID signal.}
    \label{fig:toy-plot-damped-closeup}
\end{figure}
